\method{} is an HLC-style unlearning objective that trains the unlearned model to remain forgotten after simulated inner relearning updates.
The implementation uses LoRA adapters for local GPU feasibility and simulates $K$ inner SGD steps on a relearn pool.
At selected inner times, it penalizes resurrection of forget-set items.

At a high level, each outer step uses:

\begin{enumerate}[leftmargin=*]
  \item a forget batch for negative preference or immediate forget pressure;
  \item a retain batch for utility preservation;
  \item an inner sequence of relearning batches;
  \item resurrection penalties after one or more simulated relearning steps;
  \item a top-$k$ utility KL term against the full model.
\end{enumerate}

The current strongest Qwen tuning candidate uses:
\[
K=4,\quad \eta_{\mathrm{inner}}=2\cdot 10^{-5},\quad
\eta_{\mathrm{outer}}=2\cdot 10^{-5},
\]
\[
\lambda_0=1,\quad \lambda_K=4,\quad \lambda_R=2,\quad
\lambda_{\mathrm{growth}}=1,
\]
with resurrection penalties at inner steps $1,2,4$ and margin-growth target $1$.

The method is best understood as a local diagnostic objective rather than a completed SOTA algorithm.
Its purpose is to test whether explicitly controlling post-relearning resurrection can move the conditional survival curve after immediate matching.
The current evidence shows that this objective improves some short-horizon signals, but the final Qwen seeded long-horizon comparison is still needed before making a positive claim.
